% !TeX program = latexmk --lualatex --interaction=nonstopmode --halt-on-error --output-directory=build --synctex=1

\documentclass{article}
\usepackage[a4paper, margin=30mm]{geometry}

\usepackage[
    colorlinks,
    citecolor=black,
    linkcolor=black,
]{hyperref}

\usepackage{amsmath}
\usepackage{amssymb}
\numberwithin{equation}{section}

\usepackage{cancel}
\usepackage[
    per-mode=fraction
]{siunitx}

\newcommand{\dd}{\mathrm{d}}
\newcommand{\pd}[2]{\dfrac{\partial #1}{\partial #2}}

\newcommand{\unitvec}[1]{\vec{e}_{#1}}

% \renewcommand{\div}[1]{\text{div}\left(#1\right)}
\renewcommand{\div}{\nabla \cdot}

\setlength{\parindent}{0pt}

\title{Navier--Stokes in Polar Coordinates and Rotationally Symmetric Spherical Coordinates}
\author{Gidon Bauer}

\begin{document}

%%% RWTH Colors
% Hausfarbe
\definecolor{RWTHBlau100}{RGB}{0,84,159}	% Logo primary color
\definecolor{RWTHBlau75}{RGB}{64,127,183}
\definecolor{RWTHBlau50}{RGB}{142,186,229}	% Logo secondary color
\definecolor{RWTHBlau25}{RGB}{199,221,242}
\definecolor{RWTHBlau10}{RGB}{232,241,250}

% Sekundärfarben
% Schwarz
\definecolor{RWTHSchwarz100}{RGB}{0,0,0}
\definecolor{RWTHSchwarz75}{RGB}{100,101,103}
\definecolor{RWTHSchwarz50}{RGB}{156,158,159}
\definecolor{RWTHSchwarz25}{RGB}{207,209,210}
\definecolor{RWTHSchwarz10}{RGB}{236,237,237}

% Magenta
\definecolor{RWTHMagenta100}{RGB}{227,0,102}
\definecolor{RWTHMagenta75}{RGB}{233,96,136}
\definecolor{RWTHMagenta50}{RGB}{241,158,177}
\definecolor{RWTHMagenta25}{RGB}{249,210,218}
\definecolor{RWTHMagenta10}{RGB}{253,238,240}

% Gelb
\definecolor{RWTHGelb100}{RGB}{255,237,0}
\definecolor{RWTHGelb75}{RGB}{255,240,85}
\definecolor{RWTHGelb50}{RGB}{255,245,155}
\definecolor{RWTHGelb25}{RGB}{255,250,209}
\definecolor{RWTHGelb10}{RGB}{255,253,238}

% Petrol
\definecolor{RWTHPetrol100}{RGB}{0,97,101}
\definecolor{RWTHPetrol75}{RGB}{45,127,131}
\definecolor{RWTHPetrol50}{RGB}{125,164,167}
\definecolor{RWTHPetrol25}{RGB}{191,208,209}
\definecolor{RWTHPetrol10}{RGB}{230,236,236}

% Türkis
\definecolor{RWTHTuerkis100}{RGB}{0,152,161}
\definecolor{RWTHTuerkis75}{RGB}{0,177,183}
\definecolor{RWTHTuerkis50}{RGB}{137,204,207}
\definecolor{RWTHTuerkis25}{RGB}{202,231,231}
\definecolor{RWTHTuerkis10}{RGB}{235,246,246}

% Grün
\definecolor{RWTHGruen100}{RGB}{87,171,39}
\definecolor{RWTHGruen75}{RGB}{141,192,96}
\definecolor{RWTHGruen50}{RGB}{184,214,152}
\definecolor{RWTHGruen25}{RGB}{221,235,206}
\definecolor{RWTHGruen10}{RGB}{242,247,236}

% Maigrün
\definecolor{RWTHMaigruen100}{RGB}{189,205,0}
\definecolor{RWTHMaigruen75}{RGB}{208,217,92}
\definecolor{RWTHMaigruen50}{RGB}{224,230,154}
\definecolor{RWTHMaigruen25}{RGB}{240,243,208}
\definecolor{RWTHMaigruen10}{RGB}{249,250,237}	

% Orange
\definecolor{RWTHOrange100}{RGB}{246,168,0}
\definecolor{RWTHOrange75}{RGB}{250,190,80 }
\definecolor{RWTHOrange50}{RGB}{253,212,143}
\definecolor{RWTHOrange25}{RGB}{254,234,201}
\definecolor{RWTHOrange10}{RGB}{255,247,234}	

% Rot
\definecolor{RWTHRot100}{RGB}{204,7,30}
\definecolor{RWTHRot75}{RGB}{216,92,65}
\definecolor{RWTHRot50}{RGB}{230,150,121}
\definecolor{RWTHRot25}{RGB}{243,205,187}
\definecolor{RWTHRot10}{RGB}{250,235,227}	

% Bordeaux
\definecolor{RWTHBordeaux100}{RGB}{161,16,53}
\definecolor{RWTHBordeaux75}{RGB}{182,82,86}
\definecolor{RWTHBordeaux50}{RGB}{205,139,135}
\definecolor{RWTHBordeaux25}{RGB}{229,197,192}
\definecolor{RWTHBordeaux10}{RGB}{245,232,229}

% Violett
\definecolor{RWTHViolett100}{RGB}{97,33,88}
\definecolor{RWTHViolett75}{RGB}{131,78,117}
\definecolor{RWTHViolett50}{RGB}{168,133,158}
\definecolor{RWTHViolett25}{RGB}{210,192,205}
\definecolor{RWTHViolett10}{RGB}{237,229,234}	

% Lila
\definecolor{RWTHLila100}{RGB}{122,111,172}
\definecolor{RWTHLila75}{RGB}{155,145,193}
\definecolor{RWTHLila50}{RGB}{188,181,215}
\definecolor{RWTHLila25}{RGB}{222,218,235}
\definecolor{RWTHLila10}{RGB}{242,240,247}

\maketitle

\section{Navier--Stokes Equations in Polar Coordinates}

\begin{gather}
    \pd{\vec{u}}{t} + \div \left[
        \vec{u} \otimes \vec{u} +
        \frac{p}{\rho} \mathbb{I} -
        \frac{\mu}{\rho} \left(\nabla \vec{u} + \nabla \vec{u}^{\top}\right)
    \right] = 0 \\
    \div \vec{u} = 0
\end{gather}

\begin{center}
    \def\arraystretch{2}
    \begin{tabular}[]{|c|c|}
        \hline
        $x = r \cos(\theta)$ & $r = \sqrt{x^2 + y^2}$ \\
        $y = r \sin(\theta)$ & $\theta = \arctan\left(\frac{y}{x}\right)$ \\
        \hline
        \hline
        $\pd{x}{r} = \cos(\theta)$         & $\pd{r}{x} = \cos(\theta)$ \\
        $\pd{x}{\theta} = -r \sin(\theta)$ & $\pd{r}{y} = \sin(\theta)$ \\
        $\pd{y}{r} = \sin(\theta)$         & $\pd{\theta}{x} = -\frac{\sin(\theta)}{r}$ \\
        $\pd{y}{\theta} = r \cos(\theta)$  & $\pd{\theta}{y} = \frac{\cos(\theta)}{r}$ \\
        \hline
    \end{tabular}
\end{center}

\begin{align*}
        u_x = \cos(\theta) u_r - \sin(\theta) u_\theta \\
        u_y = \sin(\theta) u_r + \cos(\theta) u_\theta \\
\end{align*}

\begin{align*}
        u_x = \cos(\theta) u_r - \sin(\theta) u_\theta \\
        u_y = \sin(\theta) u_r + \cos(\theta) u_\theta \\
\end{align*}

\begin{equation}
    \div \vec{u}
    = \pd{u_x}{x} + \pd{u_y}{y}
    = \pd{u_r}{r} + \frac1r \pd{u_\theta}{\theta} + \frac{u_r}{r}
\end{equation}

\begin{equation}
    \begin{aligned}
        \nabla \vec{u}
        &= \pd{u_x}{x} \unitvec{x} \otimes \unitvec{x} +
          \pd{u_x}{y} \unitvec{x} \otimes \unitvec{y} +
          \pd{u_y}{x} \unitvec{y} \otimes \unitvec{x} +
          \pd{u_y}{y} \unitvec{y} \otimes \unitvec{y}
        = \left(\begin{matrix}
            \pd{u_x}{x} & \pd{u_x}{y} \\
            \pd{u_y}{x} & \pd{u_y}{y} \\
        \end{matrix}\right) \\
        &=
            \left(\frac1r \pd{u_\theta}{\theta} + \frac{u_r}{r}\right) \vec{e}_\theta \otimes \vec{e}_\theta +
            \pd{u_\theta}{r} \vec{e}_\theta \otimes \unitvec{r} +
            \left(\frac1r \pd{u_r}{\theta} - \frac{u_\theta}{r}\right) \unitvec{r} \otimes \vec{e}_\theta +
            \pd{u_r}{r} \unitvec{r} \otimes \unitvec{r} \\
        &= \left(\begin{matrix}
            \dfrac1r \pd{u_\theta}{\theta} + \dfrac{u_r}{r} & \pd{u_\theta}{r} \\
            \dfrac1r \pd{u_r}{\theta} - \dfrac{u_\theta}{r} & \pd{u_r}{r} \\
        \end{matrix}\right) \\
    \end{aligned}
\end{equation}

\begin{equation}
    \begin{aligned}
        \div T &=
            \left(\pd{T_{xx}}{x} + \pd{T_{yx}}{y}\right) \unitvec{x} +
            \left(\pd{T_{xy}}{x} + \pd{T_{yy}}{y}\right) \unitvec{y} \\
        &=
            \left(\frac1r \pd{T_{\theta \theta}}{\theta} + \pd{T_{r\theta}}{r} + \frac{T_{r\theta} + T_{\theta r}}{r}\right) \vec{e}_\theta +
            \left(\frac1r \pd{T_{\theta r}}{\theta} + \pd{T_{rr}}{r} + \frac{T_{rr} - T_{\theta \theta}}{r}\right) \unitvec{r}
    \end{aligned}
\end{equation}

\begin{align*}
    T_{\theta \theta} &=
    u_\theta^2
    + \frac{p}{\rho}
    - 2\frac{\mu}{\rho} \left(
        \frac1r \pd{u_\theta}{\theta} + \frac{u_r}{r}
    \right) \\
    T_{r \theta} &= T_{\theta r} = u_r u_\theta - \frac{\mu}{\rho} \left(
        \pd{u_\theta}{r} + \frac1r \pd{u_r}{\theta} - \frac{u_\theta}{r}
    \right) \\
    T_{rr} &= u_r^2 + \frac{p}{\rho} - 2\frac{\mu}{\rho} \pd{u_r}{r}
\end{align*}

\begin{equation}
    \begin{aligned}
        \pd{u_\theta}{t} &+
        \frac1r \pd{}{\theta} \left[
            u_\theta^2 +
            \frac{p}{\rho} -
            \frac{2\mu}{\rho} \left(\dfrac1r \pd{u_\theta}{\theta} + \dfrac{u_r}{r}\right)
        \right] \\ &+
        \pd{}{r} \left[
            u_r u_\theta -
            \frac{\mu}{\rho} \left(
                \pd{u_\theta}{r} + \frac1r \pd{u_r}{\theta} - \frac{u_\theta}{r}
            \right)
        \right] \\ &+
        \frac2r \left[
            u_\theta u_r - \frac{\mu}{\rho}\left(\pd{u_\theta}{r} + \dfrac1r \pd{u_r}{\theta} - \dfrac{u_\theta}{r}\right)
        \right]
        = 0
    \end{aligned}
\end{equation}

\begin{equation}
    \pd{u_r}{r} + \frac1r \pd{u_\theta}{\theta} + \frac{u_r}{r} = 0
\end{equation}

\begin{equation}
    \begin{aligned}
        \Delta p &= \frac{\partial^2 p}{\partial x^2} + \frac{\partial^2 p}{\partial y^2} \\
        &= \frac1r \pd{}{r} \left(r \pd{p}{r}\right) + \frac{1}{r^2} \frac{\partial^2 p}{\partial \theta^2} \\
        &= \frac{\partial^2 p}{\partial r^2} +
        \frac1r \pd{p}{r} +
        \frac{1}{r^2} \frac{\partial^2 p}{\partial \theta^2}
    \end{aligned}
\end{equation}

\begin{equation}
    \pd{s}{t} + \nabla \cdot \left(s \vec{u} - D \nabla s\right) = 0
\end{equation}

\begin{equation}
    \nabla s = \pd{s}{r} \unitvec{r} + \frac1r \pd{s}{\theta} \vec{e}_\theta
\end{equation}

\subsection{Couette Flow in Polar Coordinates}

Assumptions:
\begin{itemize}
    \item $u_r = 0$
    \item $\pd{}{\theta} = 0$
    \item $\pd{}{t} = 0$
    \item $\nu := \frac{\mu}{\rho}$
\end{itemize}

\begin{equation}
    \begin{gathered}
        \pd{}{r} \left[
            -\nu \left(
                \pd{u_\theta}{r} - \frac{u_\theta}{r}
            \right)
        \right] +
        \frac2r \left[
            -\nu \left(\pd{u_\theta}{r} - \dfrac{u_\theta}{r}\right)
        \right]
        = 0 \\
        \Leftrightarrow
        -\nu \frac{\partial^2 u_\theta}{\partial r^2}
        +\nu \pd{}{r} \left(\frac{u_\theta}{r}\right)
        - \frac{2\nu}{r} \pd{u_\theta}{r}
        + \frac{2\nu}{r^2} u_\theta
        = 0 \\
        \Leftrightarrow
        -\nu \frac{\partial^2 u_\theta}{\partial r^2}
        +\frac{\nu}{r} \pd{u_\theta}{r}
        -\frac{\nu}{r^2} u_\theta
        - \frac{2\nu}{r} \pd{u_\theta}{r}
        + \frac{2\nu}{r^2} u_\theta
        = 0 \\
        \Leftrightarrow
        \frac{\partial^2 u_\theta}{\partial r^2}
        -\frac{1}{r} \pd{u_\theta}{r}
        +\frac{1}{r^2} u_\theta
        +\frac{2}{r} \pd{u_\theta}{r}
        -\frac{2}{r^2} u_\theta
        = 0 \\
        \Leftrightarrow
        \frac{\partial^2 u_\theta}{\partial r^2}
        +\frac{1}{r} \pd{u_\theta}{r}
        -\frac{1}{r^2} u_\theta
        = 0 \\
    \end{gathered}
\end{equation}

\begin{equation}
    u_\theta(r) = k_1 r + \frac{k_2}{r}
\end{equation}

Boundary conditions $u_\theta(1) = 1$ and $u_\theta(2) = 0$:
\begin{gather*}
    u_\theta(1) = k_1 + k_2 = 1 \\
    u_\theta(2) = 2 k_1 + \frac12 k_2 = 0 \\
    \Rightarrow k_1 = -\frac13, \quad k_2 = \frac43
\end{gather*}

\subsection{Channel Flow in Polar Coordinates}

Assumptions:
\begin{itemize}
    \item $u_r = 0$
    \item $\pd{}{t} = 0$
    \item $\pd{p}{\theta} = \text{const.}$
    \item $\pd{p}{r} = 0$
    \item $\nu := \frac{\mu}{\rho}$
\end{itemize}

\begin{equation}
    \div \vec{u} = \frac1r \pd{u_\theta}{\theta} = 0 \Rightarrow \pd{u_\theta}{\theta} = 0
\end{equation}

\begin{equation}
    \begin{aligned}
        &\frac1r \pd{}{\theta} \left[
            u_\theta^2 +
            \frac{p}{\rho}
        \right] +
        \pd{}{r} \left[
            -\frac{\mu}{\rho} \left(
                \pd{u_\theta}{r} - \frac{u_\theta}{r}
            \right)
        \right] +
        \frac2r \left[
            -\frac{\mu}{\rho}\left(\pd{u_\theta}{r} - \dfrac{u_\theta}{r}\right)
        \right] = 0\\
        \Leftrightarrow&\
        \frac{1}{r} \pd{p}{\theta} +
        \mu \left[
            -\frac{\partial^2 u_\theta}{\partial r^2}
            + \frac1r \pd{u_\theta}{r}
            - \frac{u_\theta}{r^2}
        \right] +
        \mu
        \left[
            - \frac2r \pd{u_\theta}{r}
            + 2 \dfrac{u_\theta}{r^2}
        \right] = 0\\
        \Leftrightarrow&\ 
        \frac{\partial^2 u_\theta}{\partial r^2}
        + \frac1r \pd{u_\theta}{r}
        - \frac{u_\theta}{r^2}
        = \frac{1}{\mu r} \pd{p}{\theta}
    \end{aligned}
\end{equation}

\begin{equation}
    u_\theta(r) = k_1r + \frac{k_2}{r} + \frac{1}{2\mu} \pd{p}{\theta} r \ln(r)
\end{equation}

\begin{equation}
    \begin{split}
        &u_\theta(1) = k_1 + k_2 = 0 \\
        &u_\theta(2) = 2k_1 + \frac{k_2}{2} + \frac{1}{\mu} \pd{p}{\theta} \ln(2) = 0 \\
        \Rightarrow &\ k_1 = k_2 = -\frac{4 \ln(2)}{5\mu} \pd{p}{\theta}
    \end{split}
\end{equation}

\section{General Orthogonal Coordinates}

\subsection{Operators 2D Orthogonal Coordinate System}

The operators in an orthogonal coordinate system $(q_1, q_2)$ in two dimensions are:
\begin{itemize}
    \item Scalar gradient:
    \begin{equation}
        \nabla f
        = \frac{\unitvec{1}}{h_1} \pd{f}{q_1}
        + \frac{\unitvec{2}}{h_2} \pd{f}{q_2}
    \end{equation}
    \item Vector divergence:
    \begin{equation}
        \nabla \cdot \vec{F} =
        \frac{1}{h_1 h_2} \left[
              \pd{h_2 F_1}{q_1}
            + \pd{h_1 F_2}{q_2}
        \right]
    \end{equation}
    \item Scalar Laplacian:
    \begin{equation}
        \Delta f = 
        \frac{1}{h_1 h_2} \left[
            \pd{}{q_1} \left(\frac{h_2}{h_1} \pd{f}{q_1}\right)
            + \pd{}{q_2} \left(\frac{h_1}{h_2} \pd{f}{q_2}\right)
        \right]
    \end{equation}
    \item Vector gradient:
    \begin{equation}
        \nabla \vec{F} = \left(\begin{matrix}
            \dfrac{1}{h_1} \pd{F_1}{q_1} + \dfrac{F_2}{h_1 h_2} \pd{h_1}{q_2} &
            \dfrac{1}{h_1} \pd{F_2}{q_1} - \dfrac{F_1}{h_1 h_2} \pd{h_1}{q_2} \\
            \dfrac{1}{h_2} \pd{F_1}{q_2} - \dfrac{F_2}{h_1 h_2} \pd{h_2}{q_1} &
            \dfrac{1}{h_2} \pd{F_2}{q_2} + \dfrac{F_1}{h_1 h_2} \pd{h_2}{q_1} \\
        \end{matrix}\right)
    \end{equation}
    \item Tensor divergence:
    \begin{equation}
        \nabla \cdot T = \left(\begin{matrix}
            \dfrac{1}{h_1 h_2} \left[
                \pd{h_2 T_{11}}{q_1}
                + \pd{h_1 T_{21}}{q_2}
                + T_{12} \pd{h_1}{q_2}
                - T_{22} \pd{h_2}{q_1}
            \right] \\
            \dfrac{1}{h_1 h_2} \left[
                \pd{h_2 T_{12}}{q_1}
                + \pd{h_1 T_{22}}{q_2}
                + T_{21} \pd{h_2}{q_1}
                - T_{11} \pd{h_1}{q_2}
            \right]
        \end{matrix}\right)
    \end{equation}
\end{itemize}

The flux term of the Navier--Stokes equation becomes
\begin{equation}
    \begin{gathered}
        \vec{u} \otimes \vec{u} 
        + \frac{p}{\rho} \mathbb{I}
        - \nu \left(\nabla \vec{u} + \nabla \vec{u}^\top\right) = \\
        \vec{u} \otimes \vec{u} 
        + \frac{p}{\rho} \mathbb{I}
        - \nu \left(\begin{matrix}
            2 \left(\dfrac{1}{h_1} \pd{u_1}{q_1} + \dfrac{u_2}{h_1 h_2} \pd{h_1}{q_2}\right) &
            \dfrac{1}{h_1} \pd{u_2}{q_1}
            + \dfrac{1}{h_2} \pd{u_1}{q_2}
            - \dfrac{u_1}{h_1 h_2} \pd{h_1}{q_2}
            - \dfrac{u_2}{h_1 h_2} \pd{h_2}{q_1} \\
            \dfrac{1}{h_1} \pd{u_2}{q_1}
            + \dfrac{1}{h_2} \pd{u_1}{q_2}
            - \dfrac{u_1}{h_1 h_2} \pd{h_1}{q_2}
            - \dfrac{u_2}{h_1 h_2} \pd{h_2}{q_1} &
            2 \left(\dfrac{1}{h_2} \pd{u_2}{q_2} + \dfrac{u_1}{h_1 h_2} \pd{h_2}{q_1}\right)
        \end{matrix}\right)
    \end{gathered}
\end{equation}


\subsection{Operators 3D Orthogonal Coordinate System}

We have an orthogonal coordinate system $(q_1, q_2, q_3)$ in three dimensions with the metric $h_1$, $h_2$, $h_3$, and $H = h_1h_2h_3$.
The operators are:
\begin{itemize}
    \item Scalar gradient:
    \begin{equation}
        \nabla f
        = \frac{\unitvec{1}}{h_1} \pd{f}{q_1}
        + \frac{\unitvec{2}}{h_2} \pd{f}{q_2}
        + \frac{\unitvec{3}}{h_3} \pd{f}{q_3}
    \end{equation}
    \item Vector divergence:
    \begin{equation}
        \nabla \cdot \vec{F} =
        \frac{1}{H} \left[
              \pd{}{q_1} \left(\frac{H}{h_1} F_1\right)
            + \pd{}{q_2} \left(\frac{H}{h_2} F_2\right)
            + \pd{}{q_3} \left(\frac{H}{h_3} F_3\right)
        \right]
    \end{equation}
    \item Scalar Laplacian:
    \begin{equation}
        \Delta f = 
        \frac{1}{H} \left[
              \pd{}{q_1} \left(\frac{H}{h_1^2} \pd{f}{q_1}\right)
            + \pd{}{q_2} \left(\frac{H}{h_2^2} \pd{f}{q_2}\right)
            + \pd{}{q_3} \left(\frac{H}{h_3^2} \pd{f}{q_3}\right)
        \right]
    \end{equation}
    \item Vector gradient:
    \begin{equation}
        \begin{split}
            \left(\nabla \vec{F}\right)_{ij} &=
            \frac{1}{h_i} \pd{F_j}{q_i} - \frac{F_i}{h_i h_j} \pd{h_i}{q_j}, \quad i \ne j \\
            \left(\nabla \vec{F}\right)_{ii} &=
            \frac{1}{h_i} \pd{F_i}{q_i} + \sum_{k \ne i} \frac{F_k}{h_i h_k} \pd{h_i}{q_k}
        \end{split}
    \end{equation}
    \begin{equation}
        \nabla \vec{u} = \left(\begin{matrix}
            \dfrac{1}{h_1} \pd{u_1}{q_1} +
            \dfrac{u_2}{h_1 h_2} \pd{h_1}{q_2} +
            \dfrac{u_3}{h_1 h_3} \pd{h_1}{q_3} & 
            \dfrac{1}{h_1} \pd{u_2}{q_1} - \dfrac{u_1}{h_1 h_2} \pd{h_1}{q_2} &
            \dfrac{1}{h_1} \pd{u_3}{q_1} - \dfrac{u_1}{h_1 h_3} \pd{h_1}{q_3} \\
            \dfrac{1}{h_2} \pd{u_1}{q_2} - \dfrac{u_2}{h_2 h_1} \pd{h_2}{q_1} &
            \dfrac{1}{h_2} \pd{u_2}{q_2} +
            \dfrac{u_1}{h_2 h_1} \pd{h_2}{q_1} +
            \dfrac{u_3}{h_2 h_3} \pd{h_2}{q_3} &
            \dfrac{1}{h_2} \pd{u_3}{q_2} - \dfrac{u_2}{h_2 h_3} \pd{h_2}{q_3} \\
            \dfrac{1}{h_3} \pd{u_1}{q_3} - \dfrac{u_3}{h_3 h_1} \pd{h_3}{q_1} &
            \dfrac{1}{h_3} \pd{u_2}{q_3} - \dfrac{u_3}{h_3 h_2} \pd{h_3}{q_2} &
            \dfrac{1}{h_3} \pd{u_3}{q_3} +
            \dfrac{u_2}{h_3 h_2} \pd{h_3}{q_2} +
            \dfrac{u_3}{h_3 h_3} \pd{h_3}{q_3}
        \end{matrix}\right)
    \end{equation}
    \item Tensor divergence:
    \begin{equation}
        T = \left(\begin{matrix}
            T_{11} & T_{12} & T_{13} \\
            T_{21} & T_{22} & T_{23} \\
            T_{31} & T_{32} & T_{33} \\
        \end{matrix}\right)
    \end{equation}
    \begin{equation}
        \left(\nabla \cdot T\right)_j =
        \frac{1}{H} \sum_i \pd{}{q_i} \left(
            \frac{H}{h_i} T_{ij}
        \right) + 
        \sum_{i \ne j} \left(
            \frac{T_{ji}}{h_ih_j} \pd{h_j}{q_i} -
            \frac{T_{ii}}{h_ih_j} \pd{h_i}{q_j}
        \right)
    \end{equation}
    \begin{equation}
        \begin{split}
            \left(\nabla \cdot T\right)_1 =
            &\frac{1}{H} \left[
                  \pd{}{q_1} \left(\frac{H}{h_1} T_{11}\right)
                + \pd{}{q_2} \left(\frac{H}{h_2} T_{21}\right)
                + \pd{}{q_3} \left(\frac{H}{h_3} T_{31}\right)
            \right] + \\
            &\left(
                \frac{T_{12}}{h_2h_1} \pd{h_1}{q_2} -
                \frac{T_{22}}{h_2h_1} \pd{h_2}{q_1}
            \right) +
            \left(
                \frac{T_{13}}{h_3h_1} \pd{h_1}{q_3} -
                \frac{T_{33}}{h_3h_1} \pd{h_3}{q_1}
            \right)
        \end{split}
    \end{equation}
\end{itemize}

The flux term of the Navier--Stokes equation becomes
\begin{equation}
    \vec{u} \otimes \vec{u} 
    + \frac{p}{\rho} \mathbb{I}
    - \nu \left(\nabla \vec{u} + \nabla \vec{u}^\top\right)
\end{equation}

\subsubsection{Symmetry in Thrid Dimension}

For symmetry in the third dimension, we assume $\frac{\partial (\cdot)}{\partial q_3} = 0$ and  $u_3 = 0$.
We then have the following operators:
\begin{itemize}
    \item Scalar gradient:
    \begin{equation}
        \nabla f
        = \frac{\unitvec{1}}{h_1} \pd{f}{q_1}
        + \frac{\unitvec{2}}{h_2} \pd{f}{q_2}
    \end{equation}
    \item Vector divergence:
    \begin{equation}
        \nabla \cdot \vec{F} =
        \frac{1}{H} \left[
              \pd{}{q_1} \left(\frac{H}{h_1} F_1\right)
            + \pd{}{q_2} \left(\frac{H}{h_2} F_2\right)
        \right]
    \end{equation}
    \item Scalar Laplacian:
    \begin{equation}
        \Delta f = 
        \frac{1}{H} \left[
              \pd{}{q_1} \left(\frac{H}{h_1^2} \pd{f}{q_1}\right)
            + \pd{}{q_2} \left(\frac{H}{h_2^2} \pd{f}{q_2}\right)
        \right]
    \end{equation}
    \item Vector gradient:
    \begin{equation}
        \begin{split}
            \left(\nabla \vec{F}\right)_{ij} &=
            \frac{1}{h_i} \pd{F_j}{q_i} - \frac{F_i}{h_i h_j} \pd{h_i}{q_j}, \quad i \ne j \\
            \left(\nabla \vec{F}\right)_{ii} &=
            \frac{1}{h_i} \pd{F_i}{q_i} + \sum_{k \ne i} \frac{F_k}{h_i h_k} \pd{h_i}{q_k}
        \end{split}
    \end{equation}
    \begin{equation}
        \nabla \vec{u} = \left(\begin{matrix}
            \dfrac{1}{h_1} \pd{u_1}{q_1} +
            \dfrac{u_2}{h_1 h_2} \pd{h_1}{q_2} & 
            \dfrac{1}{h_1} \pd{u_2}{q_1} - \dfrac{u_1}{h_1 h_2} \pd{h_1}{q_2} &
            \textcolor{RWTHRot100}{0} \\
            \dfrac{1}{h_2} \pd{u_1}{q_2} - \dfrac{u_2}{h_2 h_1} \pd{h_2}{q_1} &
            \dfrac{1}{h_2} \pd{u_2}{q_2} +
            \dfrac{u_1}{h_2 h_1} \pd{h_2}{q_1} &
            \textcolor{RWTHRot100}{0} \\
            \textcolor{RWTHRot100}{0} &
            \textcolor{RWTHRot100}{0} &
            \textcolor{RWTHRot100}{
                \dfrac{u_1}{h_3 h_1} \pd{h_3}{q_1} +
                \dfrac{u_2}{h_3 h_2} \pd{h_3}{q_2}
            }
        \end{matrix}\right)
    \end{equation}
    \item Tensor divergence:
    \begin{equation}
        T = \left(\begin{matrix}
            T_{11} & T_{12} & T_{13} \\
            T_{21} & T_{22} & T_{23} \\
            T_{31} & T_{32} & T_{33} \\
        \end{matrix}\right)
    \end{equation}
    \begin{equation}
        \left(\nabla \cdot T\right)_j =
        \frac{1}{H} \sum_i \pd{}{q_i} \left(
            \frac{H}{h_i} T_{ij}
        \right) + 
        \sum_{i \ne j} \left(
            \frac{T_{ji}}{h_ih_j} \pd{h_j}{q_i} -
            \frac{T_{ii}}{h_ih_j} \pd{h_i}{q_j}
        \right)
    \end{equation}
    \begin{equation}
        \begin{split}
            \left(\nabla \cdot T\right)_1 &=
            \frac{1}{H} \left[
                  \pd{}{q_1} \left(\frac{H}{h_1} T_{11}\right)
                + \pd{}{q_2} \left(\frac{H}{h_2} T_{21}\right)
            \right] \\
            &+ \left(
                \frac{T_{12}}{h_2h_1} \pd{h_1}{q_2} -
                \frac{T_{22}}{h_2h_1} \pd{h_2}{q_1}
            \right) \\
            &- \textcolor{RWTHRot100}{
                \frac{T_{33}}{h_3h_1} \pd{h_3}{q_1}
            }
        \end{split}
    \end{equation}
    \begin{equation}
        \begin{split}
            \left(\nabla \cdot T\right)_2 &=
            \frac{1}{H} \left[
                  \pd{}{q_1} \left(\frac{H}{h_1} T_{12}\right)
                + \pd{}{q_2} \left(\frac{H}{h_2} T_{22}\right)
            \right] \\
            &+ \left(
                \frac{T_{21}}{h_1h_2} \pd{h_2}{q_1} -
                \frac{T_{11}}{h_1h_2} \pd{h_1}{q_2}
            \right) \\
            &- \textcolor{RWTHRot100}{
                \frac{T_{33}}{h_3h_2} \pd{h_3}{q_2}
            }
        \end{split}
    \end{equation}
\end{itemize}

\subsection{Metrics for Common Coordinate Systems}

\begin{itemize}
    \item Cartesian:
    \begin{align*}
        &h_1 = h_x = 1 \\
        &h_2 = h_y = 1 \\
        &h_3 = h_z = 1
    \end{align*}
    \item Polar:
    \begin{align*}
        &h_1 = h_\theta = r \\
        &h_2 = h_r = 1 \\
        &h_3 = h_\phi = 1
    \end{align*}
    \item Spherical:
    \begin{align*}
        &h_1 = h_\theta = r \\
        &h_2 = h_r = 1 \\
        &h_3 = h_\phi = r \sin(\theta)
    \end{align*}
\end{itemize}

\end{document}